\subsection{Interface}

The three work-stealing deque implementations below share the same programming interface \refLib{zoo_saturn/ws_deque_2.mli}:

\begin{description}
\item[\ocamlinline{val create : unit -> 'a t}] returns a new deque.
\end{description}

\noindent ``Private'' operations may only be called by the domain owning the deque:
\begin{description}
\item[\ocamlinline{val size : 'a t -> int}] returns the number of elements in the deque.

\item[\ocamlinline{val is_empty : 'a t -> bool}] checks if the deque is empty.

\item[\ocamlinline{val push : 'a t -> 'a -> unit}] pushes an element to the ``back'' the deque.

\item[\ocamlinline{val pop : 'a t -> 'a option}] pops an element from the ``back'' of the deque.
\end{description}

\noindent Finally, \ocamlinline{steal} is ``public'' and therefore may be called by any domain:
\begin{description}
\item[\ocamlinline{val steal : 'a t -> 'a option}] attempts to steal
  an element from the ``front'' of the deque.
\end{description}
